\chapter{从数域到函数域}\label{L9}\lecturedate{2020/03/16}\sourcepage{72}
本讲介绍函数域的Langlands对应。
\section{作为整体域的函数域}\label{L9S1}
整数环与域$k$上的一元多项式环$k[t]$有许多相似之处。$\Z$中的素数对应于$k[t]$中的不可约多项式。如前所述，每个素数$p$可以看作$\Q$上的一个赋值；类似地，每个不可约多项式$r(t)$在分式域$k(t)$上给出赋值
\begin{equation}\tag{1.1}\label{L9eq11}
\left|r(t)^a\frac{u(t)}{v(t)}\right|_r=c^{-a},
\end{equation}
其中$c>1$固定，$u(t),v(t)\in k[t]$，且$r(t)\nmid u(t)$、$r(t)\nmid v(t)$。

对固定的$c>1$，还有赋值$|\cdot|_\infty$，定义为
\begin{equation}\tag{1.2}\label{L9eq12}
\left|\frac{u(t)}{v(t)}\right|_\infty=c^{\deg u-\deg v}.
\end{equation}
虽然$|\cdot|_\infty$不同于每个$|\cdot|_r$，但它们全都是非阿基米德的！与具有阿基米德素位\idx{阿基米德素位}的$\Q$不同，$k(t)$有如下性质。
\originalnumber{1}
\begin{lemma}\label{L9S1O1}
$k(t)$上在$k$上平凡的非平凡赋值，只可能等价于\eqref{L9eq11}或\eqref{L9eq12}中的赋值。特别地，若$k$为有限域，则$k(t)$的所有非平凡赋值都是非阿基米德的。
\end{lemma}
\sourcepage{73}
\begin{proof}
先证明有限域只有平凡赋值。若$k$有限，则对任意$0\ne a\in k$，存在$n\geq1$使$a^n=1$。因此，$1=|1|=|a^n|=|a|^n$，从而$|a|=1$。

设$|\cdot|$为$k(t)$上在$k$上平凡的非平凡赋值。存在不可约多项式$r(t)$，使$|r(t)|\ne1$。选择这样一个次数最小的$r(t)$，便有：对所有非零$d(t)\in k[t]$，若$\deg d<\deg r$，则$|d(t)|=1$。若$\deg r>1$，依假设$|t|=1$。由三角不等式，对任意$g(t)\in k[t]$，有
\[|g(t)|\leq|t^{\deg g}|+\cdots+|1|=\deg g+1.\]
令$g(t)=r(t)^n$，则
\[|r(t)|^n\leq n\deg r+1,\qquad |r(t)|\leq\lim_{n\to\infty}(n\deg r+1)^{1/n}=1.\]
故$R:=|r(t)|<1$。对$r\nmid u$，将$u(t)$写成
\[u(t)=a_m(t)r(t)^m+\cdots+a_0(t),\quad a_0(t)\ne0,\quad\deg a_j<\deg r.\]
由三角不等式，
\[|u(t)|\leq R^m+\cdots+1=\frac{1-R^{m+1}}{1-R}=\frac{1-R^{\lfloor\deg u/\deg r\rfloor+1}}{1-R}.\]
对$u(t)^n$作同样的展开，得到
\[|u(t)|^n=|u(t)^n|\leq\frac{1-R^{\lfloor n\deg u/\deg r\rfloor+1}}{1-R}\leq\frac1{1-R}.\]
令$n\to\infty$，可得$|u(t)|\leq1$。另一方面，
\[1=|a_0(t)^n|\leq|u(t)|^n+R.\]
若$|u(t)|<1$，令$n\to\infty$会得到$1\leq R<1$的矛盾。因此，凡$r\nmid u$都有$|u(t)|=1$，于是$|\cdot|$具有\eqref{L9eq11}的形式，其中$c=1/R$。
% [Tutor Note: 原书 p73 由此处 $1=|a_0(t)^n|\leq|u(t)|^n+R$ 直接断言 $|u(t)|=1$；严格的说法是：若 $|u(t)|<1$，则 $|u(t)|^n\to0$，与 $1\leq|u(t)|^n+R$（其中 $R<1$）矛盾，故 $|u(t)|\geq1$，结合前文的 $|u(t)|\leq1$ 才得到 $|u(t)|=1$。译文补出了这一步。]

若$\deg r=1$，则可写$r(t)=t+a$，$a\in k$。若$|r(t)|<1$，上述展开论证同样给出\eqref{L9eq11}。因此，只需考虑$|r(t)|>1$。对任意$a\ne b\in k$，展开$(t+b+(a-b))^n$得到
\[|t+a|^n\leq\sum_{j=0}^n|t+b|^j,\]
取$n$次方根并令$n\to\infty$，有$|t+a|\leq\max\{|t+b|,1\}$。因$|t+a|>1$，必有$|t+b|>1$，于是$|t+a|\leq|t+b|$。利用$a,b$的对称性，得到$|t+a|=c>1$与$a$无关。证明对任意非零$u(t)\in k[t]$都有$|u(t)|=c^{\deg u(t)}$的方法类似，留给读者。
\end{proof}
\sourcepage{74}
设$k=\F_q$为有限域。对不可约$r(t)$，剩余域$k[t]/(r(t))\cong\F_{q^{\deg r}}$。取$c=q^{\deg r}$时，\eqref{L9eq11}中的赋值称为\textbf{归一化赋值}；\eqref{L9eq12}中的$|\cdot|_\infty$在$c=q$时称为归一化的。

记$K_*$为$k(t)$关于$|\cdot|_*$的完备化，其中$*=r$或$\infty$。注意，$K_\infty\cong k((t^{-1}))$。记$\OO_*$为赋值环。$K_*$上存在不变测度$\mu$，称为Haar测度；它在相差一个乘法常数的意义下唯一。可以通过$\mu(\OO_*)=1$将它归一化。
\originalnumber{2}
\begin{theorem}\label{L9S1O2}
若$|\cdot|_*$为归一化赋值，$\mu$为$K_*$上的归一化Haar测度，则对任意$\alpha,\beta\in K_*$，有$\mu(\alpha+\beta\OO_*)=|\beta|_*$。
\end{theorem}
\begin{proof}
由$\mu$的不变性，只需证明$\mu(\beta\OO_*)=|\beta|_*$。以下设$\beta\ne0$。若$*=r$，写
\[\beta=\sum_{j\geq j_0}a_j(t)r(t)^j,\]
其中$a_j$或者为零，或者$\deg a_j<\deg r$，且$a_{j_0}\ne0$。于是$\beta/r^{j_0}\in\OO_r^\times$，因为$|\beta/r^{j_0}|_r=1$。若$j_0\geq0$，令$S$为次数小于$\deg r$的多项式（包括0）组成的集合，则$|S|=q^{\deg r}$，并且
\[\OO_r=\bigsqcup_{(\alpha_1,\ldots,\alpha_{j_0})\in S^{j_0}}\left(\sum_{l=1}^{j_0}\alpha_l r(t)^{l-1}+\beta\OO_r\right).\]
因此，$\mu(\beta\OO_r)q^{j_0\deg r}=\mu(\OO_r)=1$，从而$\mu(\beta\OO_r)=q^{-j_0\deg r}=|\beta|_r$。

若$j_0<0$，则
\[\beta\OO_r=\bigsqcup_{(\alpha_1,\ldots,\alpha_{-j_0})\in S^{-j_0}}\left(\sum_{l=1}^{-j_0}\alpha_l r(t)^{-l}+\OO_r\right).\]
于是$\mu(\beta\OO_r)=\mu(\OO_r)q^{-j_0\deg r}=q^{-j_0\deg r}=|\beta|_r$。

若$*=\infty$，则$\OO_\infty=\{\sum_{n\geq0}a_nt^{-n}\mid a_n\in k\}$。令$\deg\beta=m_0$，则
\[\begin{aligned}
\OO_\infty&=\bigsqcup_{(\alpha_{m_0+1},\ldots,\alpha_0)\in\F_q^{-m_0}}\left(\sum_{l=m_0+1}^{0}\alpha_l t^l+\beta\OO_\infty\right)&& (m_0\leq0),\\
\beta\OO_\infty&=\bigsqcup_{(\alpha_1,\ldots,\alpha_{m_0})\in\F_q^{m_0}}\left(\sum_{l=1}^{m_0}\alpha_l t^l+\OO_\infty\right)&& (m_0>0).
\end{aligned}\]
两种情形都给出$\mu(\beta\OO_\infty)=q^{m_0}=|\beta|_\infty$，结论得证。
\end{proof}
\sourcepage{75}
我们要研究的函数域比有理函数域$k(t)$稍复杂一些。
\originalnumber{3}
\begin{definition}{曲线的函数域}\label{L9S1O3}
设$X$为有限域$\F_q$上的光滑射影连通曲线。$X$上定义在$\F_q$上的有理函数组成的域$\F_q(X)$，称为$X$的\textbf{函数域}\znote{「函数域」与数域相对：固定有限域$\F_q$后，$\F_q$上一条光滑射影曲线$X$对应一个超越次数为$1$的域$\F_q(X)$，且由$X\mapsto\F_q(X)$可以复原$X$（曲线范畴与函数域范畴反变等价）。与$\Q$的对应关系是：素位对应$X$的闭点（即不可约多项式），整闭包$D_X$对应$\Z$，分式理想对应除子。所以函数域是「几何版的数域」，这正是把它当作整体域来研究的原因。[来源：nLab "function field"，\url{https://ncatlab.org/nlab/show/function+field}（Last revised 2022-06-05；2026-11-14 访问）。]}。
\end{definition}
\originalnumber{4}
\begin{note}[L9S1O4]
函数域$\F_q(X)$是$\F_q(t)$的有限扩张。$\F_q(t)$上的每个赋值只有有限多个延拓到$\F_q(X)$；由下面的\eref{引理}{L9S1O5}，它们都是非阿基米德的。
\end{note}
\originalnumber{5}
\begin{lemma}\label{L9S1O5}
域$K$上的赋值$|\cdot|$为非阿基米德赋值，当且仅当对$K$中由1生成的子环中的所有$n$，都有$|n|\leq1$。特别地，若$\operatorname{char}(K)>0$，则$K$的任何赋值都为非阿基米德赋值。
\end{lemma}
\begin{proof}
按定义，$|\cdot|$非阿基米德，当且仅当$|\alpha+\beta|\leq\max\{|\alpha|,|\beta|\}$。这一条件显然推出$|n|\leq1$。反方向只需证明：若所有$|n|\leq1$，则对任意$0<|\alpha|\leq1$，有$|1+\alpha|\leq1$。由三角不等式，
\[|1+\alpha|^n=|(1+\alpha)^n|\leq\sum_{j=0}^n\left|\binom nj\right||\alpha|^j\leq n+1.\]
故$|1+\alpha|\leq(1+n)^{1/n}$。令$n\to\infty$，得到$|1+\alpha|\leq1$。
\end{proof}
$\F_q[t]$在$\F_q(X)$中的整闭包$D_X$是Dedekind整环。对任意非零素理想$\mathfrak R\subset D_X$，完备化$\F_q(X)_{\mathfrak R}$是$\F_{q^m}((\widetilde t))$的有限扩张，其中$\mathfrak R\cap\F_q[t]=(r(t))$，$r(t)$不可约，$\widetilde t=r(t)$，$m=\deg r(t)$。令$N=[\F_q(X)_{\mathfrak R}:\F_{q^m}((\widetilde t))]$，则作为向量空间，$\F_q(X)_{\mathfrak R}\cong\F_{q^m}((\widetilde t))^{\oplus N}$。每个$\alpha$给出线性变换$m_\alpha:x\mapsto\alpha x$。定义
\[\operatorname{Norm}_{\F_q(X)_{\mathfrak R}/\F_{q^m}((\widetilde t))}(\alpha):=\det m_\alpha\in\F_{q^m}((\widetilde t)).\]
基域上的归一化赋值$|\cdot|_r$唯一延拓到$\F_q(X)_{\mathfrak R}$，为
\[|\alpha|=\left|\operatorname{Norm}_{\F_q(X)_{\mathfrak R}/\F_{q^m}((\widetilde t))}(\alpha)\right|_r^{1/N}.\]
在未完备的整体域情形，有如下定理。
\originalnumber{6}
\begin{theorem}\label{L9S1O6}
设$|\cdot|$为整体域$k$的某个素位\znote{「素位」（place）是数域上「素数」的推广：它指一个绝对值的等价类。它可以来自素数，给出非阿基米德绝对值（如$\Q$上的$p$进绝对值）；也可以来自实嵌入或复嵌入，给出阿基米德绝对值（如$\Q$上通常的绝对值）。用「素位」而不用「素数」，正是为了同时容纳这两类；在函数域情形所有素位都是非阿基米德的，于是素位与曲线上的闭点一一对应。[来源：nLab "place"，\url{https://ncatlab.org/nlab/show/place}（Last revised 2015-12-17；2026-11-14 访问）。]}上的归一化赋值，$K/k$为有限扩张。则
\[\prod_{j=1}^J|\alpha|_j=|\operatorname{Norm}_{K/k}(\alpha)|,\]
\sourcepage{76}
其中$|\cdot|_j$为$|\cdot|$在$K$上的各个延拓所对应的归一化赋值。
\end{theorem}
函数域$F:=\F_q(X)$作为整体域，还具有以下性质：
\begin{itemize}
\item 对任意$\alpha\in F^\times$，只有有限多个归一化赋值满足$|\alpha|>1$；
\item $\prod_*|\alpha|_*=1$，乘积遍历$F$的所有归一化赋值；
\item 相应的弱逼近定理和强逼近定理成立。
\end{itemize}
希望曲线定义在有限域上的另一个原因，是其函数域在素位处完备化后得到局部紧的局部域。
\originalnumber{7}
\begin{theorem}\label{L9S1O7}
设$K$为带有非平凡非阿基米德赋值$|\cdot|$的域。则$K$局部紧，当且仅当以下三个条件成立：
\begin{enumerate}\item $K$完备；\item 剩余域有限；\item 赋值离散。\end{enumerate}
\end{theorem}
\section{自守表示与Galois表示}\label{L9S2}
\textbf{自守表示。}

设$F=\F_q(X)$为函数域。同样可以通过限制直积定义阿代尔环\znote{阿代尔环$\A_F$把所有完备化$F_v$用「限制直积」装在一起：除有限多个分量外，其余分量都落在整数环$\OO_v$里。它比直积小得多（直积太庞大，紧性完全丧失），又比$F$大得多，因此既携带全部局部信息，又是局部紧的。伊代尔群$I_F=\A_F^\times$同理，只是要求每个分量可逆。整体域上的分析、Hecke理论都在$\A_F$上进行，这是「局部—整体原则」的技术载体。[来源：nLab "ring of adeles"，\url{https://ncatlab.org/nlab/show/ring+of+adeles}（Last revised 2025-07-01；2026-11-14 访问）。]}$\A_F$与伊代尔群$I_F$。作为拓扑环，$\A_F$局部紧。$\GL_n(\A_F)$定义为元素取自$\A_F$的可逆$n\times n$矩阵组成的群。其限制直积拓扑使$\GL_n(\A_F)$局部紧，并具有紧开
\sourcepage{77}
子群$K:=\GL_n(\prod_v\OO_v)=\prod_v\GL_n(\OO_v)$。固定$\GL_n(\A_F)$上的Haar测度，使$\mu(K)=1$。令$G=\GL_n$。
\originalnumber{1}
\begin{definition}{尖点自守函数空间}\label{L9S2O1}
设$\mathcal A_0(G,\omega)$为$G(\A_F)$上满足以下条件的所有复值函数$f$组成的空间：
\begin{enumerate}[label=(\roman*)]
\item \textbf{光滑性：}$f$在$G(\A_F)$上局部常值。
\item \textbf{$K$有限性：}$f$的所有右$K$平移张成有限维向量空间，即
\[\dim_{\C}\operatorname{span}_{\C}\{h\cdot f\mid h\in K\}<\infty.\]
\item \textbf{中心特征标：}设$Z(\A_F)$为$G(\A_F)$中的标量矩阵群，且
% [Tutor Note: 原书 p77 作 "the diagonal matrices"，但 $G(\A_F)$ 的中心应为标量矩阵而非全部对角矩阵，译文按标量矩阵处理。]
\[\omega=\prod_v\omega_v:Z(\A_F)\cong\A_F^\times\to\C^\times\]
为在$F^\times$上平凡的特征标。对每个$z\in Z(\A_F)$、$\gamma\in G(F)$，要求$f(z\gamma x)=\omega(z)f(x)$。
\item \textbf{尖点性：}对每个真标准抛物子群$P_{n_1,\ldots,n_s}=M_{n_1,\ldots,n_s}N_{n_1,\ldots,n_s}$，其中$n_1+\cdots+n_s=n$、$s\geq2$，$N$为幂幺根基，积分
\[f_{N_{n_1,\ldots,n_s}}(x):=\int_{N_{n_1,\ldots,n_s}(F)\backslash N_{n_1,\ldots,n_s}(\A_F)}f(gx)\,d\mu(g)\]
作为$G(\A_F)$上的函数恒为零。
\end{enumerate}
\end{definition}
\originalnumber{2}
\begin{note}[L9S2O2]
函数域上定义$\mathcal A(G,\omega)$所需的条件比数域情形少，因为$\A_F$没有来自阿基米德素位的分量。这使我们初步感觉，函数域情形应当比数域情形容易处理。
\end{note}
重数一定理对$\mathcal A_0(G)$同样成立。每个不可约自守表示$\pi$都可以写成限制张量积
\[\pi\cong\bigotimes_{v\in|X|}'\pi_v,\]
其中几乎所有$\pi_v$都是球面表示，即$\pi_v^{G(\OO_v)}\ne0$。在处处非分歧的情形，所有$\pi_v$都是球面的。
\sourcepage{78}
为在局部理论中分类不可约$\pi_v$，同样使用Hecke理论。赋值
\[\pi_v\longmapsto\pi_v^{G(\OO_v)}\]
在$G(F_v)$的不可约球面表示的同构类与Hecke代数$\mathcal H(\GL_n(F_v),\GL_n(\OO_v))$的不可约模之间建立双射。后者可以用$\GL_n(\C)$中的半单共轭类来描述，而$\pi_v^{G(\OO_v)}$也是一维的。

因此，在自守一侧，为每个不可约尖点表示$\pi$赋予一族$\GL_n(\C)$中的半单共轭类：
\[\pi\longleftrightarrow\{\gamma_v\}_{v\text{为非分歧素位}}.\]
这一侧的理论与数域情形类似，只是$\pi$没有来自阿基米德素位的分量。

\textbf{Galois表示。}

注意，$F$不是完美域。这里以$F^{\mathrm{sep}}$表示其可分闭包。设$v$为$F$的一个素位，亦即$X$的一个闭点。$\Gal(F^{\mathrm{sep}}/F)$中有分解子群$D_v$，它在共轭意义下唯一，且同构于$\Gal(F_v^{\mathrm{sep}}/F_v)$。映射$D_v\to\Gal(\overline{k(v)}/k(v))$的核为惯性子群$I_v$。若$I_v\subset\ker\sigma$，称Galois表示$\sigma$在$v$处非分歧。这时，$\mathrm{Fr}_v\in\Gal(\overline{k(v)}/k(v))$给出表示中的良定义共轭类。

再次选择$\overline{\Q}_\ell$向量空间来实现Galois表示，其中$\ell\ne\operatorname{char}(F)$。在处处非分歧的情形，可以把这些表示构造为$X$上的光滑$\ell$进层，亦即局部系统\znote{「局部系统」就是局部常值层：$\Q_\ell$向量空间的局部系统在每个足够小的开集上平凡，粘合数据由基本群（在$\ell$进情形是$\Gal(F^{\mathrm{sep}}/F)$）的表示给出。所以「$\sigma$在处处非分歧」相当于说这个局部系统能延拓到整条$X$上，而不只是某个开子集上。[来源：nLab "local system"，\url{https://ncatlab.org/nlab/show/local+system}（Last revised 2025-04-10；2026-11-14 访问）。]}。因此，对每个这样的$\sigma$，得到一族$\GL_n(\overline{\Q}_\ell)$中的共轭类$\widetilde\sigma_v$。令$\sigma_v$为$\widetilde\sigma_v$的半单部分，即对角矩阵$\operatorname{diag}(\lambda_1(v),\ldots,\lambda_n(v))$的共轭类，其中$\lambda_j(v)$是$\widetilde\sigma_v$的特征值。由下面的定理可见，若几乎所有$v\in|X|$都有$\sigma_{1,v}=\sigma_{2,v}$，则$\sigma_1^{\mathrm{ss}}\cong\sigma_2^{\mathrm{ss}}$；这里$\sigma_j^{\mathrm{ss}}$为$\sigma_j$的半单化。特别地，不可约表示$\sigma$由这族$\{\sigma_v\}$所确定。
\originalnumber{3}
\begin{theorem}{Chebotarev密度定理}\label{L9S2O3}\idx{Chebotarev密度定理}
设$L/F$为整体域的有限Galois扩张，$C\subset\Gal(L/F)$为共轭类。则存在无限多个非分歧素位$\p$，使$\mathrm{Fr}_{\p}$的共轭类为$C$。
\end{theorem}
\sourcepage{79}
\textbf{Langlands对应。}

函数域情形已经是一条定理：$n=2$时由Drinfeld证明，$n>2$时由Lafforgue证明。
\originalnumber{4}
\begin{theorem}\label{L9S2O4}
存在自守表示与Galois表示之间的双射，使
\[\pi\longleftrightarrow\sigma,\qquad\sigma_v=\gamma_v\in\GL_n(\overline{\Q}),\quad\forall v\in|X|.\]
\znote{原书未写出下列条件；以下为文献中的标准表述，供对照，不属原书内容：需固定$\ell\ne\operatorname{char}(F)$及系数域之间的同构$\overline{\Q}_\ell\cong\C$（这两项约定本讲正文已在讨论$\ell$进表示处说明，此处不重复）；$\pi$应取具有有限阶中心特征标的不可约尖点自守表示，$\sigma$应取行列式具有有限阶的不可约连续$n$维$\ell$进Galois表示，且后者几乎处处非分歧。此即Drinfeld（$n=2$）与Lafforgue（$n>2$）的定理。[来源：本讲参考文献 [D] Drinfeld, Amer. J. Math. 105 (1983) 85-114；[L] Lafforgue, Invent. Math. 147 (2002) 1-241。]}
\end{theorem}
\needspace{5\baselineskip}\section{第9讲参考文献}\label{L9S3}
\textbf{一般性介绍。}
\begin{references}
\refitem{L9bibBG}{BG} J. Bernstein and S. Gelbart, eds, \emph{An introduction to the Langlands Program}, Birkhäuser, 2004, Chapter 12 written by D. Gaitsgory.
\end{references}
\textbf{技术细节。}
\begin{references}
\refitem{L9bibD}{D} V. G. Drinfeld, \emph{Two-dimensional $\ell$-adic representations of the fundamental group of a curve over a finite field and automorphic forms on GL(2)}, Amer. J. Math. 105 (1983) 85-114.
% [Tutor Note: 原书 p79 此条作者名印作 "V. G. Drinfel"，应为 V. G. Drinfeld，译文按正确拼写处理。]
\refitem{L9bibL}{L} L. Lafforgue, \emph{Chtoucas de Drinfeld et correspondance de Langlands}, Invent. Math. 147 (2002) 1-241.
\end{references}
\needspace{4\baselineskip}\textbf{关于Chebotarev密度定理。}
\begin{references}
\refitem{L9bibFJ}{FJ} Michael D. Fried and Moshe Jarden, \emph{Field Arithmetic}, Chapter 7. The Chebotarev Density Theorem, Springer-Verlag Berlin, Heidelberg 2008.
\end{references}
